%% 自動生成（gen.py）
\bktocchap{2-0}{「だんだん\R{速}{はや}く」を\R{数}{かず}で\R{言}{い}うと？}{ことばを\R{数}{かず}と\R{単位}{たん|い}にすると、くらべられて\R{計算}{けい|さん}もできる}{1}{1}{1}{8}{J}
\bktocpart{\R{第}{だい}1\R{部}{ぶ}}{\R{量}{りょう}と\R{数}{かず}}{16}
\bktocchap{2-1}{「かぞえる\R{数}{かず}」と「はかる\R{数}{かず}」}{「5」だけでは、\R{何}{なん}のことかわからない}{7}{1}{1}{18}{C}
\bktocchap{2-2}{マイナスの\R{数}{かず}と\R{数直線}{すう|ちょく|せん}}{0\,℃より\R{寒}{さむ}いときは、どう\R{書}{か}く？}{6}{1}{2}{38}{C}
\bktocchap{2-3}{＋ － × ÷ は、\R{何}{なに}を\R{表}{あらわ}しているの？}{かけ\R{算}{ざん}の\R{答}{こた}えは、\R{長方形}{ちょう|ほう|けい}の\R{広}{ひろ}さ？}{1,2}{1}{1}{60}{C}
\bktocchap{2-4}{\R{分数}{ぶん|すう}と\R{小数}{しょう|すう}の\R{計算}{けい|さん}}{ピザ2\R{枚}{まい}を3\R{人}{にん}で\R{分}{わ}けると、1\R{人分}{にん|ぶん}は\R{何枚}{なん|まい}？}{1}{1}{2}{76}{C}
\bktocchap{2-5}{\R{比}{ひ}・\R{割合}{わり|あい}・\R{百分率}{ひゃく|ぶん|りつ}}{「100\,\%」をこえることはある？}{7}{1}{2}{104}{C}
\bktocchap{2-6}{はかった\R{数}{かず}には「ずれ」がある}{\R{同}{おな}じ\R{物}{もの}を10\R{回}{かい}はかると、ぜんぶ\R{同}{おな}じになる？}{7}{1}{2}{122}{C}
\bktocpart{\R{第}{だい}2\R{部}{ぶ}}{\R{大}{おお}きな\R{数}{かず}・\R{小}{ちい}さな\R{数}{かず}・\R{単位}{たん|い}}{148}
\bktocchap{2-7}{\R{何回}{なん|かい}もかける：\R{累乗}{るい|じょう}と\R{指数}{し|すう}}{\R{紙}{かみ}を42\R{回}{かい}\R{折}{お}ると、\R{月}{つき}まで\R{届}{とど}く？}{3}{1}{2}{150}{C}
\bktocchap{2-8}{2\R{回}{かい}かけると□になる\R{数}{かず}：\R{平方根}{へい|ほう|こん}}{\R{面積}{めん|せき}が2\,m²の\R{正方形}{せい|ほう|けい}、1\R{辺}{ぺん}は\R{何}{なん}\,m？}{3}{2}{2}{168}{C}
\bktocchap{2-9}{\R{分数}{ぶん|すう}で\R{書}{か}けない\R{数}{かず}がある？}{$\sqrt{2}$は、どこまで\R{計算}{けい|さん}しても\R{終}{お}わらない}{6}{2}{2}{184}{C}
\bktocchap{2-10}{\R{単位}{たん|い}の\R{世界}{せ|かい}ルール：SI\R{単位}{たん|い}}{\R{時速}{じ|そく}72\,kmは、\R{秒速}{びょう|そく}\R{何}{なん}\,m？}{7}{1}{2}{202}{C}
\bktocpart{\R{第}{だい}3\R{部}{ぶ}}{\R{文字}{も|じ}と\R{式}{しき}}{228}
\bktocchap{2-11}{なぜ\R{文字}{も|じ}を\R{使}{つか}うの？}{どんな\R{速}{はや}さでも\R{使}{つか}える\R{式}{しき}を、1\R{行}{ぎょう}で}{5}{2}{2}{230}{C}
\bktocchap{2-12}{\R{文字式}{も|じ|しき}のやくそく}{3\,mと2\R{秒}{びょう}は、たし\R{算}{ざん}できない}{5}{2}{2}{256}{C}
\bktocchap{2-13}{\R{公式}{こう|しき}に\R{数}{かず}を\R{入}{い}れてみよう：\R{代入}{だい|にゅう}}{\R{式}{しき}に\R{数}{かず}を\R{入}{い}れると、\R{答}{こた}えが\R{出}{で}てくる}{5}{2}{2}{272}{C}
\bktocchap{2-14}{\R{式}{しき}を\R{変形}{へん|けい}して、ほしい\R{量}{りょう}をとり\R{出}{だ}そう}{1つの\R{公式}{こう|しき}から、3つの\R{公式}{こう|しき}が\R{作}{つく}れる}{5}{2}{2}{294}{C}
\bktocchap{2-15}{\R{比例}{ひ|れい}と\R{反比例}{はん|ぴ|れい}}{ばねを2\R{倍}{ばい}のばすには、2\R{倍}{ばい}の\R{力}{ちから}}{1}{1}{2}{314}{C}
\bktocchap{2-16}{\R{何秒後}{なん|びょう|ご}に\R{追}{お}いつく？1\R{次方程式}{じ|ほう|てい|しき}}{\R{先}{さき}に\R{出}{で}た\R{弟}{おとうと}に、\R{兄}{あに}は\R{何秒}{なん|びょう}で\R{追}{お}いつく？}{1}{2}{2}{334}{C}
\bktocpart{\R{第}{だい}4\R{部}{ぶ}}{コンピュータで\R{計算}{けい|さん}}{360}
\bktocchap{2-17}{コンピュータの\R{数}{かず}：2\R{進数}{しん|すう}と16\R{進数}{しん|すう}}{0と1だけで、どんな\R{数}{かず}も\R{書}{か}ける}{3}{2}{2}{362}{C}
\bktocchap{2-18}{\R{表計算}{ひょう|けい|さん}で\R{計算}{けい|さん}しよう}{1000\R{回}{かい}の\R{計算}{けい|さん}も、コピー1\R{回}{かい}で}{5}{2}{2}{386}{C}
\bktocchap{2-19}{プログラムで\R{計算}{けい|さん}しよう}{0.1を10\R{回}{かい}たしても、1にならない？}{1}{2}{2}{406}{C}
\par\vspace{2mm}
\bktocbacklist{\bkbi{こたえ}{426}\bkbi{\R{記号}{き|ごう}カード}{438}\bkbi{よく\R{使}{つか}う\R{定数}{てい|すう}・\R{単位}{たん|い}・\R{公式}{こう|しき}}{441}\bkbi{できるようになったことチェック\R{表}{ひょう}}{446}\bkbi{つぎの\R{巻}{かん}へ}{449}\bkbi{さくいん}{452}}
